\documentclass[10pt,a4paper]{amsart}
\usepackage[margin=19mm]{geometry}
\usepackage{amsmath,amssymb,mathtools}
\usepackage[scr=rsfs,cal=euler]{mathalfa}
\usepackage{xfrac}
\usepackage[unicode,hidelinks,bookmarks=false]{hyperref}
\newcommand{\ie}{i.e.\ }
\newcommand{\cf}{cf.\ }
\newcommand{\resp}{respectively\ }
\newcommand{\Subsection}{\subsection}
\providecommand{\nobreakdash}{\nobreak\hbox{-}}
\setlength{\emergencystretch}{2em}
\multlinegap0pt
\usepackage{bm}
\usepackage{bbm}
\usepackage{dsfont}
\usepackage{xspace}
\usepackage{cancel}
\usepackage{mathtools}
\usepackage{amsthm}
\makeatletter
\@ifundefined{proposition}{\newtheorem{proposition}{Proposition}}{}
\makeatother
\title[IGA Laplace Eigenfrequencies Distributions and Estimations: ]{IGA Laplace Eigenfrequencies Distributions and Estimations: Impact of Reparametrization on Eigenfrequency Behavior}
\author{Lamsahel Noureddine, Abdeladim El Akri, Ahmed Ratnani}
\date{Source: arXiv:2510.12632v1 (CC BY 4.0)}
\begin{document}
\maketitle

\noindent\emph{Source and licence.} IGA Laplace Eigenfrequencies Distributions and Estimations: Impact of Reparametrization on Eigenfrequency Behavior. Version arXiv:2510.12632v1 (CC BY 4.0), distributed under the Creative Commons Attribution 4.0 International licence, as recorded in the arXiv licence metadata for this exact version and shown on the abstract page https://arxiv.org/abs/2510.12632v1. Full rights notice: licenses/arxiv-2510.12632-CC-BY-4.0.txt.\par\smallskip

\noindent\textit{Editorial scope.} Counted unit: the Proposition whose source label is \texttt{propc} in the article (source0.tex lines 488--490, source label \texttt{propc}), delivered together with the article's own complete original proof of it (source0.tex lines 493--584, one \texttt{proof} environment of 92 lines). No mathematics is written, repaired or re-derived: statement and proof are the article's own text line for line. Required context retained uncounted: largpsidif (lines 447--449); symboldif (lines 451--453); stefa (lines 462--472). Every definition, display and label of those lines is retained unchanged. Omitted from this package and not redistributed: 1--446, 450--450, 454--461, 473--487, 491--492, 585--1341. Nothing inside a retained excerpt is omitted or altered: every retained body is delivered exactly as the article's lines are written, so no omission receipt of any kind accompanies this package. Citations: the retained excerpts carry no citation command at all, so no bibliography entry of the article is transcribed and no reference is redistributed. Reference adaptation: none was needed and none was made; every reference inside the delivered unit resolves inside it, so no unresolved reference or citation is produced. Printed slips kept literal: no printed word, symbol, hypothesis or number of the counted statement or of its proof was altered. Layout and class: the article's own class is \texttt{article} with packages including geometry, graphicx, caption, subcaption, hyperref, amsmath, amsthm, amssymb, authblk, bm, xcolor, float, booktabs, yhmath; the wrapper is the lane's pinned \texttt{amsart} editorial wrapper at 10pt on A4, which restates the theorem-like environments the retained text needs and adds the article's own macro definitions that the retained text needs (none), with extra wrapper packages (bm, bbm, dsfont, xspace, cancel, mathtools). The delivered numbering is the unit's own. Assets: no retained span contains a figure environment, a graphics-inclusion command, a PDF-page inclusion, a table or a TikZ picture; no image file is copied into this package, the package declares no asset dependency and no image file is redistributed with it. Licence: version arXiv:2510.12632v1 of this article is distributed under the Creative Commons Attribution 4.0 International licence, as recorded in the arXiv licence metadata for this exact version and shown on the abstract page https://arxiv.org/abs/2510.12632v1. A full rights notice is delivered at licenses/arxiv-2510.12632-CC-BY-4.0.txt. Characters: every retained excerpt is pure ASCII, so the delivered engine, XeLaTeX with its default Latin Modern text font, reports no missing character.
\begin{equation}\label{largpsidif}
    \Psi_{\phi}^{p}(y)=\mu_{2} \left( \left\{  (x,\theta)\in [0,1]\times[0,\pi]:\;\; \sqrt{\omega_{\phi}^p(x,\theta)}\leq y    \right\} \right),\;\;\forall y\in  R_{\sqrt{\omega_{\phi}^p}},
\end{equation}

\begin{equation}\label{symboldif}
  \sqrt{\xi_{\phi}^p}(x)=\inf\left\{ y\in   R_{\sqrt{\omega_{\phi}^p}}:\;\; \Psi_{\phi}^{p}(y)>\pi  x  \right\},\;\;\forall x\in(0,1).  
\end{equation}

\begin{proposition}\label{stefa}
For all  $\phi\in \mathbf{C}_{[0,1]}$ and for all $p\geq 1$, we have 
$$
\Psi_{\phi}^{p}\in \mathcal{C}^1((0,+\infty))\;\;\text{and} \;\; \sup_{y \geq 0} \left(\Psi^{p}_{\phi}\right)'(y)<\infty.
$$ 
Moreover,
\begin{equation}\label{relatioxiandPsi}
   \sqrt{\xi_{\phi}^p}(x)=\left(\Psi_{\phi}^{p}\right)^{-1}(\pi x),\;\;\forall x \in [0,1],
\end{equation}
where $\Psi_{\phi}^{p}$ and $\sqrt{\xi_{\phi}^p}$ are given by \eqref{largpsidif} and \eqref{symboldif}, respectively.
\end{proposition}

\begin{proposition}\label{propc}
Let $\phi$ be a strictly convex reparametrization of the interval $[0,1]$. Then for every $\epsilon\in(0,1)$, the function $\Psi_{\phi}^{1}$ is $C^{2}\left(\left(0,\epsilon\frac{\sqrt{12}}{\phi^{'}(1)}  \right) \right)$ and \\ $\left(\Psi_{\phi}^{1}\right)^{'}(0)=1$. In addition, $\Psi_{\phi}^{1}$ is strictly concave over $\left(0,\frac{\sqrt{6}}{\phi^{'}(1)}\right)$.
\end{proposition}

\begin{proof}
Let $\phi\in\mathbf{C}_{[0,1]}$ to be strictly convex and $y\in \mathbb{R}_+$. As a first step, we obtain an explicit expression for $\Psi_{\phi}^{1}(y)$. We proceed as follows
\begin{align}\label{eq:propc-eq1}
%\vspace{-1cm}
\Psi_{\phi}^{1}(y)&= \mu_2\left\{ \sqrt{\omega_{\phi}^1}\leq y\right\} \nonumber\\
& =\mu_2\left\{ (x,\theta)\in [0,1]\times[0,\pi],\;\;\; \omega_{\phi}^1(x,\theta)\leq y^2 \right\} \nonumber \\
 & =\mu_2\left\{ (x,\theta)\in [0,1]\times[0,\pi],\;\;\; \dfrac{6(1-\cos{\theta})}{2+\cos{\theta}}\leq \left(y \phi^{'}(x)\right)^2 \right\} \nonumber\\
 &= \mu_2\left\{ (x,\theta)\in [0,1]\times[0,\pi],\;\;\; 6(1-\cos{\theta})\leq 2 \left(y \phi^{'}(x)\right)^2 + \left(y \phi^{'}(x)\right)^2 \cos{\theta} \right\}\nonumber\\
 &= \mu_2\left\{ (x,\theta)\in [0,1]\times[0,\pi],\;\;\; \cos{\theta}\geq \beta_x \right\}.
% \vspace{-0.54cm}
\end{align}
%\vspace{-2cm}
Here, $\beta_x$ is defined as
$$
\beta_x = \dfrac{6-2 \left(y\phi^{'}(x)\right)^2}{6+\left(y \phi^{'}(x)\right)^2}, \quad 0 \leq x \leq 1. 
%\vspace{-0.1cm}
$$
 %\vspace{-2cm}
To evaluate the last measure \eqref{eq:propc-eq1}, we employ the property that $\cos:\, [0, \pi] \longrightarrow [-1,1]$ is invertible, $\arccos$ being its inverse function. This requires characterizing the conditions under which $\beta_x$ is in $[-1, 1]$. In fact, it is easy to see that $\beta_x \leq 1$ for all $x \in [0,1]$, and $\beta_x \geq -1$ if, and only if
%\vspace{-1cm}
\begin{equation*}
%\vspace{-0.1cm}
    \left\{
\begin{array}{ll}
0 \leq x \leq 1 \\
0 \leq y \leq \frac{\sqrt{12}}{\phi'(1)}, 
 \end{array}
\right. \text{or} \quad \left\{
\begin{array}{ll}
0 \leq x \leq \left(\phi' \right)^{-1} \left( \frac{\sqrt{12}}{y} \right)\\
\frac{\sqrt{12}}{\phi'(1)} \leq y \leq \frac{\sqrt{12}}{\phi'(0)}.
\end{array}
\right.
%\vspace{-0.4cm}
\end{equation*}
Therefore, from \eqref{eq:propc-eq1}, we deduce that
\begin{equation}\label{psiequa}
  \Psi_{\phi}^{1}(y)=\left\{
\begin{array}{ll}
  \displaystyle\int_{0}^{1}\arccos\left(  \dfrac{6-2\left( y\phi'(x)\right)^2}{6+\left( y\phi'(x)\right)^2}\right)\;dx, &y\in J_1,\\\\

\displaystyle\int_{0}^{\left( \phi'\right)^{-1}\left(\dfrac{\sqrt{12}}{y}\right)}\arccos\left(  \dfrac{6-2\left( y\phi'(x)\right)^2}{6+\left( y\phi'(x)\right)^2}\right)\;dx \\ \hspace{4cm}+\pi\left(1-\left( \phi'\right)^{-1}\left(\dfrac{\sqrt{12}}{y}\right)\right),\\ 
&y\in J_2,
 \end{array}
\right.
\end{equation}
where $J_1:= \left(0,\dfrac{\sqrt{12}}{\phi'(1)}\right)$ and $J_2:=  \left(\dfrac{\sqrt{12}}{\phi'(1)},\dfrac{\sqrt{12}}{\phi'(0)}\right)$. Since $\left(\Psi_{\phi}^{1}\right)'$ is $C^1((0,+\infty[)$ by Proposition \ref{stefa}, we then compute $\left(\Psi_{\phi}^{1}\right)'$ on $J_1$. 

Let $y \in J_1$, we define
$$  
\varphi(x,y)= \arccos\left(  \dfrac{6-2\left( y\phi'(x)\right)^2}{6+\left( y\phi'(x)\right)^2}\right),\;\; \forall x \in [0,1],\quad y \in J_1.
$$
For almost every $x \in [0,1]$ and all $y \in J_1$, we have
\begin{align*}
\dfrac{\partial \varphi}{\partial y}(x,y)&=\dfrac{-1}{\sqrt{1-\left[ \dfrac{6-2\left( y\phi'(x)\right)^2}{6+\left( y\phi'(x)\right)^2}\right]^2}}\;\dfrac{-36 y \left(\phi'(x)\right)^2}{\left[ 6+\left( y\phi'(x)\right)^2\right]^2}\vspace*{0.25cm}\\
&=\dfrac{36 y \left(\phi'(x)\right)^2}{\sqrt{ \left(12-\left( y\phi'(x)\right)^2\right) \left[ 3\left( y\phi'(x)\right)^2\right]    }}\;\dfrac{1}{ 6+\left( y\phi'(x)\right)^2}\vspace*{0.25cm}\\
&=\dfrac{\phi'(x)}{\sqrt{1-\left(\dfrac{y\phi'(x)}{\sqrt{12}}\right)^2}}\; \dfrac{6}{6+\left( y\phi'(x)\right)^2} \vspace*{0.25cm}\\
& \leq \dfrac{\phi'(x)}{\sqrt{1-\left(\dfrac{\phi'(x)}{\phi'(1)}\right)^2}}  = \dfrac{\phi'(x) \sqrt{\phi'(1)}}{\sqrt{1+\dfrac{\phi'(x)}{\phi'(1)}}} \;\dfrac{1}{\sqrt{\phi'(1)-\phi'(x)}} \vspace*{0.25cm}\\
& \leq \frac{\left(\phi'(1)\right)^{3/2}}{\sqrt{1+\frac{\phi'(0)}{\phi'(1)}}} \;\dfrac{1}{\sqrt{\phi'(1)-\phi'(x)}},
\end{align*}
and
$$
\int_0^1 \dfrac{dx}{\sqrt{\phi'(1)-\phi'(x)}}    \leq \frac{1}{\displaystyle\inf_{[0,1]} \left\{\phi''\right\}} \int_{\phi'(0)}^{\phi'(1)} \dfrac{dx}{\sqrt{\phi'(1)-x}}  < \infty.
$$
Thus, by the Lebesgue Dominated Convergence Theorem, $\Psi_{\phi}^{1} \in C^1(J_1)$, and
\begin{equation}\label{PsiJ1}
\left(\Psi_{\phi}^{1}\right)'(y)=\displaystyle\int_0^1 \dfrac{\partial \varphi}{\partial y}(x,y)\;dx,\quad \forall y\in J_1,\quad \left(\Psi_{\phi}^{1}\right)'(0)=1.
\end{equation}

For the third step, we fix $\epsilon$  in $(0,1)$ and we aim to demonstrate  that $\left(\Psi_{\phi}^{1}\right)'$ is of class $C^{1}\left(J_1^{\epsilon}\right)$, where $$J_1^{\epsilon}=\left(0,\epsilon\frac{\sqrt{12}}{\phi^{'}(1)}  \right).$$

Let $y\in J_1^{\epsilon}$. From (\ref{PsiJ1}), we observe that for every $x\in [0,1]$ the function $\frac{\partial\varphi}{\partial y}$ is differentiable with respect to $y$ and we have 
\begin{align*}
    \frac{\partial^2\varphi}{\partial^2 y}(x,y)&=-6\phi^{'}(x)\dfrac{\frac{-\left(\frac{\phi^{'}(x)}{\sqrt{12}}\right)^2 2y}{2\sqrt{1-\left(\frac{y\phi^{'}(x)}{\sqrt{12}}\right)^2}}\left(6+\left(y\phi^{'}(x)\right)^2\right)+\sqrt{1-\left(\frac{y\phi^{'}(x)}{\sqrt{12}}\right)^2}\;2y\left(\phi^{'}(x)\right)^2}{\left(1-\left(\frac{y\phi^{'}(x)}{\sqrt{12}}\right)^2\right)\left(6+\left(y\phi^{'}(x)\right)^2\right)^2}\\
    &=-6y\left(\phi^{'}(x)\right)^3\dfrac{-\frac{1}{12}\left(6+\left(y\phi^{'}(x)\right)^2\right)+2\left(1-\left(\frac{y\phi^{'}(x)}{\sqrt{12}}\right)^2\right)}{\left(1-\left(\frac{y\phi^{'}(x)}{\sqrt{12}}\right)^2\right)^{\frac{3}{2}}\left(6+\left(y\phi^{'}(x)\right)^2\right)^2}\\
    &=\dfrac{-9y\left(\phi^{'}(x)\right)^3 \left(1-2\left(\frac{y\phi^{'}(x)}{\sqrt{12}}\right)^2\right)}{\left(1-\left(\frac{y\phi^{'}(x)}{\sqrt{12}}\right)^2\right)^{\frac{3}{2}}\left(6+\left(y\phi^{'}(x)\right)^2\right)^2}.
\end{align*}
Observe that $\frac{\partial^2\varphi}{\partial^2 y}$ is continuous over $[0,1]\times J_1^{\epsilon} $, then  $\Psi_{\phi}^{1}\in C^2(J_1^{\epsilon})$, with
\begin{equation}\label{PsipJ1}
\left(\Psi_{\phi}^{1}\right)^{''}(y)=\displaystyle\int_0^1 \dfrac{\partial^2 \varphi}{\partial^2 y}(x,y)\;dx,\quad\forall y\in J_1^{\epsilon}.
\end{equation}

To conclude the proof, let $y\in \left(0,\frac{\sqrt{6}}{\phi^{'}(1)}\right)$. Then, there exists an $\epsilon$ in $(0,1)$ such that 
\begin{align*}
    y&\leq \frac{\sqrt{6}}{\phi^{'}(1)}< \epsilon\frac{\sqrt{12}}{\phi^{'}(1)}\\
    &\Rightarrow y\in J_1^{\epsilon}\;\;\text{ and}\;\; 1-2\left(\frac{y\phi^{'}(x)}{\sqrt{12}}\right)^2\geq 1-\left(\frac{\phi^{'}(x)}{\phi^{'}(1)}\right)^2\geq 0.
\end{align*}
Hence, using the (\ref{PsipJ1}), we obtain 
$$\left(\Psi_{\phi}^{1}\right)^{''}(y)<0,\;\;\forall y\in \left(0,\frac{\sqrt{6}}{\phi^{'}(1)}\right). $$

Finally, the function $\Psi_{\phi}^{1}$ is strictly concave over $\left(0,\frac{\sqrt{6}}{\phi^{'}(1)}\right)$, which ends the proof.
\end{proof}

\end{document}
